\documentclass[a4paper,11pt]{article}
\usepackage{fontspec}
\setmainfont{Arial}
\usepackage{xspace}
\usepackage{url}
\usepackage{parskip}

% Styling and margins
\RequirePackage{color}
\RequirePackage{graphicx}
\usepackage[usenames,dvipsnames]{xcolor}
\usepackage[left=1.1cm,right=1.1cm,top=0.8cm,bottom=0.5cm]{geometry}

% Tables
\usepackage{tabularx}
\newcolumntype{C}{>{\centering\arraybackslash}X}

% Lists formatting
\usepackage{enumitem}
\setlist[itemize]{leftmargin=*, itemsep=3pt, topsep=3pt, parsep=0pt, partopsep=0pt}

% Titles format
\usepackage{titlesec}
\titleformat{\section}{\large\scshape\raggedright}{}{0em}{}[\titlerule]
\titlespacing{\section}{0pt}{8pt}{5pt}

% Hyperlinks
\usepackage[unicode, draft=false]{hyperref}
\definecolor{linkcolour}{rgb}{0,0.2,0.6}
\hypersetup{colorlinks,breaklinks,urlcolor=linkcolour,linkcolor=linkcolour}

% Icons
\usepackage{fontawesome5}

\begin{document}
\pagestyle{empty}

% Header
\begin{tabularx}{\linewidth}{@{} C @{}}
\LARGE{\textbf{Ляра Максим Витальевич}} \\[2pt]
\large{\textbf{ML Engineer / Data Scientist $\mid$ Опыт: более 2,5 лет}} \\[3pt]
\href{https://t.me/maxlyara1}{\raisebox{-0.05\height}\faTelegram \ @maxlyara1} \ $|$ \
\href{mailto:maks.lyara@gmail.com}{\raisebox{-0.05\height}\faEnvelope \ maks.lyara@gmail.com} \ $|$ \
\href{tel:+79119600866}{\raisebox{-0.05\height}\faMobile \ +7 (911) 960 08 66} \ $|$ \
\href{https://drive.google.com/drive/folders/1r9MWmdKT1crSFLui4-bHhgw5ErIWPB44?usp=sharing}{\raisebox{-0.05\height}\faCertificate \ Сертификаты} \ $|$ \
\href{https://github.com/maxlyara1}{\raisebox{-0.05\height}\faGithub \ GitHub}
\end{tabularx}
\vspace{3mm}

% Summary
\section{Обо мне}
Специализируюсь на полном цикле разработки ML-решений: от сбора бизнес-требований и развертывания (Docker, Kubernetes) до поддержки в промышленной эксплуатации. Работал с прогнозированием временных рядов, оптимизацией цепей поставок (MIP), RAG-системами и компьютерным зрением. Работаю на стыке моделирования и инженерии: довожу модели до рабочих сервисов и оцениваю результат по бизнес-показателям.

% Experience
\section{Опыт работы}

\noindent\textbf{Data Scientist / ML Engineer, АО «Акрихин»} \hfill \textbf{Март 2025 -- н.в.}
\vspace{3pt}
\begin{itemize}
\item \textbf{Контрактная кампания (S\&OP):} Разработал систему прогнозирования B2B-продаж (CatBoost) для аптечных сетей. Автоматизировал запуск моделей и регулярный расчёт прогнозов в Airflow.
\item \textbf{Оптимизация запасов (Safety Stocks):} Спроектировал MIP-модель (PuLP) для расчета страховых запасов. Учет бизнес-ограничений позволил снизить замороженный капитал на 10\% при поддержании уровня сервиса 98\%.
\item \textbf{LLM-агенты и RAG:} Расширял функционал корпоративной LLM-платформы (Langflow, FastAPI). Создал RAG-агентов для поиска по нормативной документации с учетом версионирования. Реализовал интеграцию с внутренними сервисами (Yandex Tracker, EQMS) с поддержкой потоковой передачи ответов (SSE).
\item \textbf{Data Mining:} Инициировал и разработал масштабируемую архитектуру конкурентного парсинга с обходом блокировок и кэшированием (TTL) для ежедневного мониторинга рынка.
\item \textbf{Аналитика графов (Антифрод):} Создал алгоритм выявления «серого реэкспорта» по графу B2B-транзакций. Использовал обходы DFS/BFS в ClickHouse для отслеживания аномальных товарных цепочек.
\item \textbf{Аудит выкладки (CV):} Автоматизировал аудит полок для фармпредставителей. Проводил дообучение (fine-tuning) YOLO-моделей для детекции и классификации широкого ассортимента SKU в аптеках.
\end{itemize}

\vspace{6pt}
\noindent\textbf{ML Engineer, Mars L.L.C.} \hfill \textbf{Декабрь 2023 -- Декабрь 2024}
\vspace{3pt}
\begin{itemize}
\item \textbf{Прогнозирование B2B-оттока:} Разработал модель прогнозирования оттока торговых точек по гео-данным (CatBoost), сократив потери дистрибуции на 7\%.
\item \textbf{CV \& NLP автоматизация:} Разработал CV-модель контроля доли полки (Share of Shelf), улучшив качество выкладки на 20\%. Внедрил NLP-модель для анализа неструктурированных текстов, ускорив обработку данных и сократив ручную разметку.
\end{itemize}

\section{Личные проекты}
\begin{itemize}
    \item \textbf{\href{https://github.com/maxlyara1/multimodal-video-search}{Multimodal Video Search} -- поиск по видео:} Система поиска по речи, тексту на экране и описаниям кадров; Qdrant, FastAPI и локальные модели. На 50 проверочных запросах лучший вариант достиг Hit@1 90\% и Hit@3 96\%.
    \item \textbf{\href{https://github.com/maxlyara1/AMML-ranking-competition}{AMML Ranking} -- ранжирование товаров:} Градиентный бустинг, признаки соответствия запроса товару и нейросетевое переранжирование результатов поиска.
    \item \textbf{\href{https://github.com/maxlyara1/find_anomalies_hackathon}{RedLab Hack} -- поиск аномалий:} Детектор аномалий в телеметрии промышленных датчиков (ETNA, Streamlit); 3-е место.
\end{itemize}

\section{Образование}
\textbf{РАНХиГС, «Цифровая экономика»} \hfill 2022 --- 2026 \\
Институт экономики, математики и информационных технологий.

\section{Навыки}
\begin{itemize}
    \item \textbf{ML \& DS:} Python, CatBoost, PyTorch, SQL, Time Series, Computer Vision (YOLO), NLP (BERT), Graphs
    \item \textbf{LLMs:} RAG (Qdrant, TEI, комбинированный поиск), ИИ-ассистенты (Langflow, FastAPI), YandexGPT, Qwen, SSE
    \item \textbf{MLOps \& Data:} Airflow, ClickHouse, PostgreSQL, Kafka, Kubernetes, Yandex DataSphere, MinIO, Docker, Git
\end{itemize}
\enlargethispage{2\baselineskip}
\end{document}
